\documentclass[../../main.tex]{subfiles}

\begin{document}

\chapter[LR Parsing]{LR Parsing: Canonical LR, SLR, LALR,\\ dan Konflik}
\label{chap:lr-parsing}

\begin{outcomebox}
\textbf{Sub-CPMK:} \subcpmk{4.2}\\[3pt]
Setelah menyelesaikan bab ini, mahasiswa mampu:
\begin{enumerate}[label=\arabic*.,leftmargin=*]
  \item menjelaskan item LR(0), closure, dan goto;
  \item membangun automaton LR(0) dan tabel ACTION/GOTO SLR(1);
  \item membangun tabel ACTION/GOTO LALR(1) dari state LR(1) yang
        digabung;
  \item menguji dan men-\textit{debug} konflik shift-reduce dan
        reduce-reduce;
  \item memilih strategi parsing LL(k) atau LR(k) sesuai kebutuhan.
\end{enumerate}
\end{outcomebox}

\section*{Materi Pokok}
\begin{itemize}[leftmargin=*]
  \item Item LR(0), closure, dan goto.
  \item Automaton LR(0) dan tabel ACTION/GOTO SLR(1).
  \item Canonical LR(1), penggabungan inti, dan tabel LALR(1).
  \item Konflik shift-reduce dan reduce-reduce.
  \item Alur debug konflik dan perbandingan LL(k)/LR(k).
\end{itemize}

\subfile{section-01}
\subfile{section-02}
\subfile{section-03}
\subfile{section-04}
\subfile{section-05}

\section*{Contoh}
\addcontentsline{toc}{section}{Contoh}

\begin{contoh}[Menghitung closure]
Diberikan item \([E \to \cdot T]\) dan aturan
\(T \to T * F \mid F\). Closure menambahkan
\([T \to \cdot T * F]\) dan \([T \to \cdot F]\). Karena \(F\)
adalah nonterminal dengan aturan \(F \to (E) \mid \code{id}\),
closure juga menambahkan \([F \to \cdot (E)]\) dan
\([F \to \cdot \code{id}]\).
\end{contoh}

\begin{contoh}[Menghitung goto]
Untuk state yang memuat \([E \to E \cdot + T]\), goto pada simbol
\(+\) menghasilkan item \([E \to E + \cdot T]\), lalu closure
menambahkan aturan untuk \(T\). State baru ini kemudian memiliki
transisi pada simbol-simbol yang dapat mengikuti \(T\).
\end{contoh}

\begin{contoh}[Mengisi tabel SLR]
State yang memuat \([E \to E + T \cdot]\) menghasilkan aksi reduce
\(E \to E + T\) pada semua token di FOLLOW(\(E\)), yaitu
\(\{\$, ), +\}\). State yang memuat \([T \to T \cdot * F]\)
menghasilkan aksi shift pada \(*\).
\end{contoh}

\begin{contoh}[Konflik dangling else]
Pada state setelah \code{if E then stmt}, token \code{else} dapat
memicu shift (memasangkan else dengan if dalam) atau reduce
(mengurangi if luar). Tanpa aturan prioritas, tabel memiliki dua
aksi pada sel yang sama. Bison biasanya memilih shift, yang sesuai
dengan aturan else terdekat.
\end{contoh}

\begin{contoh}[Penggabungan LALR]
Dua state LR(1) dengan inti yang sama, misalnya
\(\{[A \to \alpha \cdot, a]\}\) dan
\(\{[A \to \alpha \cdot, b]\}\), digabung menjadi satu state LALR
\(\{[A \to \alpha \cdot, \{a,b\}]\}\). Penggabungan mengurangi
jumlah state, tetapi dapat menimbulkan konflik reduce-reduce yang
tidak ada pada LR(1) kanonik.
\end{contoh}

\section*{Aktivitas Pembelajaran dan Praktikum}
\addcontentsline{toc}{section}{Aktivitas Pembelajaran dan Praktikum}

\begin{table}[H]
\centering
\caption{Rencana aktivitas Bab XII}
\label{tab:ch12-aktivitas}
\begin{tabularx}{\textwidth}{@{}llX@{}}
\toprule
\textbf{Sesi} & \textbf{Bentuk} & \textbf{Kegiatan} \\
\midrule
1 & Ceramah & Item LR(0), closure, dan goto \\
2 & Latihan & Membangun automaton LR(0) dan tabel SLR \\
3 & Lokakarya & Penggabungan inti LR(1) dan tabel LALR \\
4 & Praktikum & P3: tabel ACTION/GOTO dan debug konflik \\
\bottomrule
\end{tabularx}
\end{table}

\subsection*{P3 / Praktikum 3: Tabel SLR/LALR dan Debug Konflik}
\textbf{Tujuan.} Membangun tabel ACTION/GOTO untuk SLR dan LALR,
membandingkannya dengan alat, serta men-\textit{debug} konflik
shift-reduce dan reduce-reduce.

\textbf{Langkah kerja.}
\begin{enumerate}[label=\arabic*.,leftmargin=*]
  \item Bangun koleksi kanonik LR(0) untuk
        \file{spec/grammar-lr.bnf}.
  \item Isi tabel ACTION/GOTO SLR: satu baris per state, kolom ACTION
        untuk terminal, dan kolom GOTO untuk nonterminal.
  \item Bangun koleksi LR(1), gabungkan state dengan inti sama, lalu
        isi tabel ACTION dan GOTO untuk LALR.
  \item Jalankan \file{slr\_table.py} dan Bison; bandingkan jumlah
        state, isi item, dan tabel ACTION/GOTO.
  \item Identifikasi konflik shift-reduce dan reduce-reduce; catat
        state dan token pemicunya.
  \item Debug satu konflik pada satu waktu: precedence,
        restrukturisasi grammar, atau resolusi eksplisit.
  \item Uji ulang string valid, string tidak valid, dan kasus
        dangling else.
  \item Kumpulkan tabel, log sebelum/sesudah, dan laporan singkat.
\end{enumerate}

\textbf{Definisi selesai.}
\begin{checklist}
  \item Tabel ACTION/GOTO SLR lengkap.
  \item Tabel ACTION/GOTO LALR lengkap.
  \item Tidak ada konflik tak terduga; konflik yang tersisa memiliki
        resolusi eksplisit.
  \item String valid diterima dan string tidak valid ditolak.
  \item Log sebelum/sesudah dan laporan diserahkan.
\end{checklist}

Hasil praktikum menjadi bagian Praktikum 3 (12\%) yang dilaporkan
pada Bab XIII.

\section*{Latihan dan Refleksi}
\addcontentsline{toc}{section}{Latihan dan Refleksi}

\subsection*{Latihan}
\begin{enumerate}[label=\arabic*.,leftmargin=*]
  \item Definisikan item LR(0) dan berikan contoh.
  \item Hitung closure dari \([E \to \cdot T]\).
  \item Hitung goto dari state yang memuat
        \([E \to E \cdot + T]\) pada simbol \(+\).
  \item Bangun koleksi kanonik LR(0) untuk grammar
        \(S \to a S b \mid \varepsilon\).
  \item Bangun tabel SLR untuk grammar pada Latihan 4.
  \item Jelaskan perbedaan SLR, LR(1), dan LALR.
  \item Berikan contoh konflik shift-reduce.
  \item Berikan contoh konflik reduce-reduce.
\end{enumerate}

\subsection*{Refleksi}
\begin{reflectionbox}
\begin{enumerate}[label=\arabic*.,leftmargin=*]
  \item Apakah Anda dapat menjelaskan mengapa LR lebih kuat daripada
        SLR?
  \item Bagian mana dari konstruksi tabel yang paling menantang?
  \item Bagaimana Anda akan membaca laporan konflik dari Bison?
\end{enumerate}
\end{reflectionbox}

\section*{Asesmen}
\addcontentsline{toc}{section}{Asesmen}

Asesmen Bab XII mencakup Praktikum 3 (12\%) dan UTS Praktik (14\%).
Fokusnya adalah tabel ACTION/GOTO SLR dan LALR, serta kemampuan
menguji dan men-\textit{debug} konflik.

\begin{table}[H]
\centering
\caption{Instrumen asesmen Bab XII}
\label{tab:ch12-asesmen}
\begin{tabularx}{\textwidth}{@{}lXl@{}}
\toprule
\textbf{Instrumen} & \textbf{Indikator} & \textbf{Bobot} \\
\midrule
Praktikum 3 & Tabel ACTION/GOTO SLR dan LALR, log debug konflik, dan
uji string & 12\% (bagian) \\
UTS Praktik & Tabel SLR/LALR dan penyelesaian konflik shift-reduce
atau reduce-reduce & 14\% (bagian) \\
Kuis Konsep & Item LR, closure, goto, dan jenis konflik & formatif \\
\bottomrule
\end{tabularx}
\end{table}

Kriteria kelulusan: mahasiswa dapat membangun tabel ACTION/GOTO SLR
dan LALR, membaca laporan konflik, men-\textit{debug} konflik sampai
selesai, serta memilih strategi parsing yang sesuai.

\section*{Checklist Pencapaian Kompetensi}
\addcontentsline{toc}{section}{Checklist Pencapaian Kompetensi}

\begin{checklist}
  \item Saya dapat mendefinisikan item LR(0).
  \item Saya dapat menghitung closure.
  \item Saya dapat menghitung goto.
  \item Saya dapat membangun koleksi kanonik LR(0).
  \item Saya dapat membangun tabel SLR(1).
  \item Saya dapat menjelaskan perbedaan SLR, LR(1), dan LALR.
  \item Saya dapat menganalisis konflik.
  \item Saya dapat memilih strategi LL atau LR.
  \item Saya sudah menyelesaikan lokakarya verifikasi tabel.
\end{checklist}

\section*{Rangkuman}
\addcontentsline{toc}{section}{Rangkuman}

Pada bab ini telah dibahas:
\begin{itemize}[leftmargin=*]
  \item item LR(0), closure, dan goto;
  \item automaton LR(0) dan tabel SLR(1);
  \item canonical LR(1) dan LALR(1);
  \item konflik shift-reduce dan reduce-reduce;
  \item perbandingan LL(k) dan LR(k).
\end{itemize}

Bab XIII akan memakai parser generator Bison untuk membangun parser
Subset-C, menambahkan semantic action, dan menghasilkan AST. Kita
akan melihat bagaimana teori LR diterapkan dalam alat nyata.

\section*{Evaluasi}
\addcontentsline{toc}{section}{Evaluasi}

\begin{enumerate}[label=\arabic*.,leftmargin=*]
  \item Definisikan item LR(0) dan jelaskan arti posisi titik.
  \item Hitung closure dari \([A \to \cdot B C]\) dengan aturan
        \(B \to b\).
  \item Hitung goto dari state yang memuat \([A \to a \cdot B]\)
        pada simbol \(B\).
  \item Bangun tabel ACTION/GOTO SLR untuk grammar
        \(S \to A\), \(A \to a A \mid b\).
  \item Bangun tabel ACTION/GOTO LALR untuk grammar pada Soal 4.
  \item Jelaskan perbedaan SLR dan LALR.
  \item Berikan contoh konflik shift-reduce, jelaskan penyebabnya, dan
        tunjukkan cara menyelesaikannya.
  \item Berikan contoh konflik reduce-reduce, jelaskan penyebabnya, dan
        tunjukkan cara menyelesaikannya.
  \item Kapan Anda memilih LL(k) daripada LR(k)?
\end{enumerate}

\subsection*{Kunci dan Rubrik Singkat}
\begin{itemize}[leftmargin=*]
  \item Soal 1--3: ketepatan definisi dan perhitungan (70\%),
        penjelasan (30\%).
  \item Soal 4--5: kebenaran tabel (50\%), kelengkapan ACTION/GOTO
        (50\%).
  \item Soal 6--9: pemahaman konsep (60\%), argumentasi (40\%).
\end{itemize}


\ifSubfilesClassLoaded{%
  \bibliographystyle{apalike}%
  \bibliography{../../references}%
}{}

\end{document}
